%!TEX program = xelatex
\documentclass[12pt,a4paper]{exam-zh}

\usepackage{xeCJK}
\setCJKmainfont{Noto Serif CJK SC}
\usepackage{geometry}
\usepackage{amsmath,array,booktabs,tabularx,multirow,enumitem,graphicx,tikz,fancyhdr}
\geometry{margin=2cm,top=1.8cm,bottom=1.8cm}
\pagestyle{fancy}
\fancyhf{}
\cfoot{-\thepage-}
\renewcommand{\headrulewidth}{0pt}
\setlength{\parindent}{0pt}
\setlength{\parskip}{0.35em}
\linespread{1.12}

\newcommand{\optline}[4]{%
  \begin{tabular*}{\linewidth}{@{\extracolsep{\fill}}llll}
  A\;#1 & B\;#2 & C\;#3 & D\;#4
  \end{tabular*}}
\newcommand{\papertitleone}{2018年度马来西亚华文独中统一考试}
\newcommand{\papersubject}{高级数学\,(SC05)}
\newcommand{\formulasection}[1]{\vspace{0.6em}{\bfseries #1}\par}

\begin{document}

% Page 1
\begin{center}
  {\zihao{3}\bfseries \papertitleone\par}
  \vspace{1.3em}
  {\zihao{3}\bfseries 高中组\par}
  \vspace{1.2em}
  {\zihao{2}\bfseries \papersubject\par}
  \vspace{1.2em}
  {\zihao{4}\bfseries 试卷一\quad 选择题\par}
  \vspace{0.8em}
  日期：2018年10月25日\par
  时间：8.30 a.m.－9.30 a.m.\par
  （1小时）
\end{center}

\vspace{1.5em}
{\bfseries 考生须知}
\begin{enumerate}[label=（\arabic*）,leftmargin=2.6em,itemsep=0.7em]
  \item 本科试卷共分两份：试卷一：选择题（40\%）；试卷二：作答题（60\%）。
  \item 考生须于第一阶段规定的1小时内完成试卷一。暂停15分钟后，才在第二阶段规定的2小时内作答试卷二。
  \item 试卷一选择题20题全答。选出正确的答案，然后用5B铅笔将“O”答案纸（电脑卡）上相应的字母所在的小方格涂黑。
  \item 可使用所规定的电子计算机进行演算。
  \item 数学公式表在第2--4面。
  \item 试卷一选择题共印七面，最后一面空白。
\end{enumerate}

\vfill
\begin{center}
  \fbox{\rule{0pt}{1.1em}\hspace{1em}{\bfseries 未经正式宣布\quad 不得翻看内页}\hspace{1em}}
\end{center}

\vspace{1.5em}
\begin{tabular*}{\linewidth}{@{\extracolsep{\fill}}ll}
考生姓名：\underline{\hspace{5cm}} & 考生编号：\underline{\hspace{4cm}}
\end{tabular*}

\newpage


% Page 5
\begin{center}
  {\zihao{4}\bfseries 试卷一\quad 选择题}
\end{center}
\begin{enumerate}[leftmargin=2.2em,label=\arabic*.]
\item 已知 $f(x)=\dfrac{x+1}{x-4}$，求函数 $f$ 的定义域。\\
\optline{$\{x\in\mathbb R\mid x\ge 0\}$}{$\{x\in\mathbb R\mid x\le -1\text{ 或 }x>4\}$}{$\mathbb R\setminus\{4\}$}{$\{x\in\mathbb R\mid x\ge 0,\ x\ne 4\}$}

\item 化简 $\dfrac{x^4-y^4}{x^3y+xy^3}$。\\
\optline{$\dfrac{x}{y}+\dfrac{y}{x}$}{$\dfrac{x}{y}-\dfrac{y}{x}$}{$\dfrac{x}{y}+\dfrac{x}{y}$}{$\dfrac{x}{y}-\dfrac{x}{y}$}

\item 已知多项式 $p(x)$ 除以 $(x+1)(x+2)$ 所得余式是 $2x+3$，求 $p(x)$ 除以 $x+2$ 的余数。\\
\optline{$-1$}{$1$}{$2$}{$3$}

\item 求 $\displaystyle \sum_{k=1}^{100} k(k+2)$。\\
\optline{$338\,350$}{$338\,550$}{$348\,450$}{$676\,700$}

\item 已知矩阵 $\begin{pmatrix}1&2\\2&3\\0&8\end{pmatrix}$ 的逆矩阵对应元素如题图所示，求 $b$ 的值。\\
\optline{$-5$}{$-1$}{$1$}{$5$}

\item 求 $(2x-1)^9$ 的展开式中的常数项。\\
\optline{$-4032$}{$-672$}{$672$}{$4032$}

\item 求无穷等比级数 $(1+\tfrac12-\tfrac14+\cdots)$ 的和。\\
\optline{$1$}{$\dfrac43$}{$\dfrac53$}{$4$}

\item 已知一组数据 $x_1,x_2,\ldots,x_6$，且 $\sum_{i=1}^6 x_i=216$，$\sum_{i=1}^6 x_i^2=22776$，求这组数据的标准差。\\
\optline{$25.75$}{$50$}{$63.09$}{$633.06$}
\end{enumerate}

\newpage

% Page 6
\begin{enumerate}[leftmargin=2.2em,label=\arabic*.,start=9]
\item 已知 $Z$ 呈标准常态分配且 $P(Z>a)=0.7$，求 $P(Z>-a)$。\\
\optline{$0.3$}{$0.35$}{$0.5$}{$0.7$}

\item 用 \texttt{MATHEMATICS} 一字中的全部字母作排列，如果相同的字母必须相邻，有几种不同的排法？\\
\optline{$39\,916\,800$}{$4\,989\,600$}{$40\,320$}{$5\,040$}

\item 已知中国郑州市的经纬度是 $(34.8^\circ\mathrm N,116.6^\circ\mathrm E)$，日本大阪的经纬度是 $(34.8^\circ\mathrm N,135.6^\circ\mathrm E)$。求这两座城市沿着纬线的距离。\\
\optline{$1140$ 海里}{$936.1$ 海里}{$650.6$ 海里}{$570$ 海里}

\item 图1中，圆 $O$ 的半径为 $2\text{ cm}$，$A$ 与 $B$ 为圆上两点且 $\angle AOB=\dfrac{2\pi}{3}$。求阴影区域的面积。\\
\optline{$(\pi-2)\text{ cm}^2$}{$\dfrac12\text{ cm}^2$}{$(4\pi-3\sqrt3)\text{ cm}^2$}{$(4\pi-4\sqrt3)\text{ cm}^2$}

\item 以下何者是 $y=2\cos 2x$ 的图像？\\
\optline{A 图}{$B$ 图}{$C$ 图}{$D$ 图}

\begin{center}
\begin{tikzpicture}[scale=0.9]
\draw (0,0) circle (1.3);
\draw (0,0) -- (90:1.3) node[above] {$A$};
\draw (0,0) -- (-30:1.3) node[right] {$B$};
\node at (0,0) [below left] {$O$};
\fill[gray!25] (0,0) -- (90:1.3) arc[start angle=90,end angle=-30,radius=1.3] -- cycle;
\node at (2.8,-1.2) {图1};
\end{tikzpicture}
\end{center}
\end{enumerate}

\newpage

% Page 7
\begin{enumerate}[leftmargin=2.2em,label=\arabic*.,start=14]
\item 若 $7\cos\theta+\sin\theta=5$，求 $\cos 2\theta$ 的值。\\
\optline{$-\dfrac7{25}$ 或 $\dfrac7{25}$}{$\dfrac7{25}$}{$\dfrac{24}{25}$}{$-\dfrac{24}{25}$ 或 $\dfrac{24}{25}$}

\item 已知原点 $O$ 落在线段 $AB$ 上且 $AO:OB=2:3$。若点 $B$ 的坐标是 $(4,-1)$，求点 $A$ 的坐标。\\
\optline{$\left(-\dfrac83,\dfrac23\right)$}{$\left(-\dfrac43,\dfrac23\right)$}{$\left(-\dfrac23,\dfrac43\right)$}{$\left(-\dfrac23,\dfrac83\right)$}

\item 已知直线 $2ax+by=\cdots$ 的斜率为 $-\dfrac12$ 且通过点 $(-4,1)$，求 $a$ 的值。\\
\optline{$-2$}{$-1$}{$1$}{$2$}

\item 求曲线 $x^2-3xy+y^2-6=0$ 在点 $(2,1)$ 处的切线斜率。\\
\optline{$-\dfrac12$}{$-\dfrac18$}{$\dfrac18$}{$2$}

\item 已知 $\displaystyle \int (\sin 2x-2\sin x)^2\,dx = k+\cdots$，求 $k$ 的值。\\
\optline{$-4$}{$-2$}{$2$}{$4$}

\item 已知 $\displaystyle \int_0^1 f(x)\,dx=2$，$\displaystyle \int_0^2 f(x)\,dx=10$，求 $\displaystyle \int_0^1 f(2x)\,dx$。\\
\optline{$1$}{$2$}{$5$}{$10$}

\item 图2中的阴影部分是由曲线 $y=x^3-3x^2+6x-4$ 及 $y=3x$ 所围成，求其面积。\\
\optline{$\dfrac13$}{$\dfrac25$}{$\dfrac12$}{$\dfrac35$}
\end{enumerate}

\vfill
\begin{flushright}图2：原卷附图\end{flushright}

\newpage

% Page 8
\begin{center}
  {\zihao{3}\bfseries \papertitleone\par}
  \vspace{1.2em}
  {\zihao{3}\bfseries 高中组\par}
  \vspace{1.1em}
  {\zihao{2}\bfseries \papersubject\par}
  \vspace{1.2em}
  {\zihao{4}\bfseries 试卷二\quad 作答题\par}
  \vspace{0.8em}
  日期：2018年10月25日\par
  时间：9.45 a.m.－11.45 a.m.\par
  （2小时）
\end{center}

\vspace{1.2em}
{\bfseries 考生须知}
\begin{enumerate}[label=（\arabic*）,leftmargin=2.6em,itemsep=0.5em]
  \item 本科试卷共分两份：试卷一：选择题（40\%）；试卷二：作答题（60\%）。
  \item 试卷二共分两组：甲组（必答题）5题全答（20\%）；乙组（选答题）7题选答4题，但不能超过4题（40\%）。两组全部共答9题完卷。
  \item 每题必须用新的一张纸作答。
  \item 只可用蓝色或黑色原子笔书写，惟可用铅笔画图。
  \item 不必抄题，惟试题号码必须书写清楚。
  \item 所有演算必须清楚地写出。必要的几何图形必须画出。
  \item 可使用所规定的电子计算机进行演算，除非题目限制。
  \item 作答前，须在积分表（电脑卡）上正确填写考生编号及用5B铅笔将对应的小方格涂黑。作答后，须根据所答题数在试题号码栏上将有关小方格涂黑，并指出作答总数。
  \item 答卷必须依试题号码次序排列，并且将积分表（电脑卡）置于上面，合订成一本。
  \item 数学公式表在第2--4面。
  \item 试卷二作答题共印七面，最后一面空白。
\end{enumerate}

\vfill
\begin{center}\fbox{\rule{0pt}{1.1em}\hspace{1em}{\bfseries 未经正式宣布\quad 不得翻看内页}\hspace{1em}}\end{center}

\vspace{1em}
\begin{tabular*}{\linewidth}{@{\extracolsep{\fill}}ll}
考生姓名：\underline{\hspace{5cm}} & 考生编号：\underline{\hspace{4cm}}
\end{tabular*}

\newpage

% Page 9
\begin{center}{\bfseries 数学公式表}\end{center}
\formulasection{I. 代数}
\begin{align*}
ax^2+bx+c&=0\Rightarrow x=\frac{-b\pm\sqrt{b^2-4ac}}{2a}, \qquad
(a\pm b)^2=a^2\pm2ab+b^2, \\
(a\pm b)^3&=a^3\pm3a^2b+3ab^2\pm b^3, \qquad
 a^2-b^2=(a+b)(a-b), \\
a^3\pm b^3&=(a\pm b)(a^2\mp ab+b^2), \qquad
 a^m a^n=a^{m+n}, \qquad \frac{a^m}{a^n}=a^{m-n}, \\
(a^m)^n&=a^{mn}, \quad (ab)^n=a^n b^n, \quad \left(\frac ab\right)^n=\frac{a^n}{b^n}, \\
\log_a(xy)&=\log_a x+\log_a y, \quad \log_a\frac{x}{y}=\log_a x-\log_a y, \quad \log_a(x^m)=m\log_a x.
\end{align*}
\begin{align*}
a^{\log_a x}=x,\qquad \log_b x=\frac{\log_a x}{\log_a b},\qquad (a+b)^n=\sum_{r=0}^{n}{^nC_r}a^{n-r}b^r.
\end{align*}
等差数列与等比数列公式同第2面。

\formulasection{II. 三角学}
\begin{align*}
\text{弧长 }&=r\theta,\qquad \text{扇形面积 }=\frac12r^2\theta,\qquad \tan\theta=\frac{\sin\theta}{\cos\theta}, \\
\csc\theta&=\frac1{\sin\theta},\qquad \sec\theta=\frac1{\cos\theta},\qquad \cot\theta=\frac1{\tan\theta}, \\
\frac{a}{\sin A}&=\frac{b}{\sin B}=\frac{c}{\sin C}=2R,\qquad a^2=b^2+c^2-2bc\cos A, \\
\triangle\text{面积}&=\frac12ab\sin C.
\end{align*}
\begin{align*}
\sin^2\theta+\cos^2\theta&=1, & 1+\tan^2\theta&=\sec^2\theta, & 1+\cot^2\theta&=\csc^2\theta,\\
\sin(A\pm B)&=\sin A\cos B\pm\cos A\sin B,\\
\cos(A\pm B)&=\cos A\cos B\mp\sin A\sin B,\\
\tan(A\pm B)&=\frac{\tan A\pm\tan B}{1\mp\tan A\tan B},\\
\sin2A&=2\sin A\cos A,\qquad \cos2A=\cos^2A-\sin^2A=2\cos^2A-1=1-2\sin^2A,\\
\tan2A&=\frac{2\tan A}{1-\tan^2A}.
\end{align*}

\newpage

% Page 10
\begin{center}{\bfseries 数学公式表}\end{center}
\formulasection{III. 解析几何}
\begin{align*}
d&=\sqrt{(x_2-x_1)^2+(y_2-y_1)^2}, \\
\text{分比公式 }&\left(\frac{mx_2+nx_1}{m+n},\frac{my_2+ny_1}{m+n}\right), \\
\triangle\text{面积 }&=\frac12\left|x_1y_2+x_2y_3+x_3y_1-x_2y_1-x_3y_2-x_1y_3\right|, \\
\text{直线方程式 }&y-y_1=m(x-x_1), \\
\text{点到直线距离 }&=\frac{|Ax_0+By_0+C|}{\sqrt{A^2+B^2}}, \\
\text{圆的标准式 }&(x-h)^2+(y-k)^2=r^2.
\end{align*}

\formulasection{IV. 统计与概率}
\begin{align*}
\bar x&=\frac{\sum f_i x_i}{\sum f_i},\qquad \operatorname{Mo}=L+\frac{d_1}{d_1+d_2}C, \\
M&=L+\frac{\frac n2-F_m}{f_m}C,\qquad Q_3=L_3+\frac{\frac{3n}{4}-F_3}{f_3}C,\qquad Q_1=L_1+\frac{\frac n4-F_1}{f_1}C, \\
\text{四分位距 }&=Q_3-Q_1,\qquad \text{四分位差 }=\frac{Q_3-Q_1}{2}, \\
\sigma^2&=\frac{\sum f_i x_i^2}{\sum f_i}-\left(\frac{\sum f_i x_i}{\sum f_i}\right)^2,\qquad
\sigma=\sqrt{\frac{\sum f_i x_i^2}{\sum f_i}-\left(\frac{\sum f_i x_i}{\sum f_i}\right)^2}, \\
\text{平均差 }&=\frac{\sum f_i|x_i-\bar x|}{\sum f_i},\qquad \nu=\frac{\sigma}{\bar x}\times100\%, \\
r&=\frac{\sum xy-\frac{(\sum x)(\sum y)}{n}}{\sqrt{\left(\sum x^2-\frac{(\sum x)^2}{n}\right)\left(\sum y^2-\frac{(\sum y)^2}{n}\right)}}, \\
QI&=\frac{Q_1}{Q_0}\times100,\qquad \bar x_w=\frac{\sum wx}{\sum w}, \\
{}^nP_r&=\frac{n!}{(n-r)!},\qquad {}^nC_r=\frac{n!}{(n-r)!r!}, \\
P(A\cup B)&=P(A)+P(B)-P(A\cap B),\qquad P(A')=1-P(A), \\
E(X)&=x_1p_1+x_2p_2+\cdots+x_kp_k.
\end{align*}

\newpage

% Page 11
\begin{center}{\bfseries 数学公式表}\end{center}
\formulasection{V. 微积分}
\begin{align*}
\lim_{x\to0}\frac{\sin x}{x}&=1, \qquad \lim_{x\to\infty}\left(1+\frac1x\right)^x=e, \\
\frac{d}{dx}(uv)&=u\frac{dv}{dx}+v\frac{du}{dx},\qquad \frac{d}{dx}\left(\frac uv\right)=\frac{v\,du/dx-u\,dv/dx}{v^2}, \\
\frac{dy}{dx}&=\frac{dy}{du}\frac{du}{dx},\qquad \frac{d}{dx}f(g(x))=f'(g(x))g'(x), \\
\frac{d}{dx}x^n&=nx^{n-1},\qquad \int x^n\,dx=\frac{x^{n+1}}{n+1}+C\ (n\ne -1), \\
\frac{d}{dx}(\sin x)&=\cos x,\qquad \int \cos x\,dx=\sin x+C, \\
\frac{d}{dx}(\cos x)&=-\sin x,\qquad \int \sin x\,dx=-\cos x+C, \\
\frac{d}{dx}(\tan x)&=\sec^2x,\qquad \int \sec^2x\,dx=\tan x+C, \\
\frac{d}{dx}(\cot x)&=-\csc^2x,\qquad \int \csc^2x\,dx=-\cot x+C, \\
\frac{d}{dx}(\sec x)&=\sec x\tan x,\qquad \int \sec x\tan x\,dx=\sec x+C, \\
\frac{d}{dx}(\csc x)&=-\csc x\cot x,\qquad \int \csc x\cot x\,dx=-\csc x+C, \\
\frac{d}{dx}(\ln x)&=\frac1x,\qquad \frac{d}{dx}(\log_a x)=\frac{1}{x\ln a},\qquad \int \frac{1}{x}\,dx=\ln x+C, \\
\frac{d}{dx}(e^x)&=e^x,\qquad \int e^x\,dx=e^x+C,\qquad \frac{d}{dx}(a^x)=a^x\ln a,\qquad \int a^x\,dx=\frac{a^x}{\ln a}+C.
\end{align*}
面积：$\displaystyle \int_a^b y\,dx$ 或 $\displaystyle \int_c^d x\,dy$。\qquad
体积：$\displaystyle \pi\int_a^b y^2\,dx$ 或 $\displaystyle \pi\int_c^d x^2\,dy$。

\newpage

% Page 12
{\bfseries 甲组\quad 必答题\,(20\%)}\par
（本组5题全答。）

\begin{enumerate}[leftmargin=2.2em,label=\arabic*.]
\item 已知函数 $f(x)=ax+b+5$ 且 $f(-4)=-1$，$f(-1)=2$。求：
  \begin{enumerate}[label=(\alph*),leftmargin=2em]
    \item $a$ 及 $b$ 的值；\hfill (2\%)
    \item $f^{-1}(x)$。\hfill (2\%)
  \end{enumerate}

\item 图1所示为以 $O$ 为圆心的两个扇形，$O,A,C$ 及 $O,B,D$ 为一直线，$OA:OC=1:2$。若 $OA=1\text{ cm}$，扇形 $OAB$ 的周长为 $4\text{ cm}$，求阴影部分的面积。\hfill (4\%)

\item 已知圆 $x^2+y^2-4x+2y-4=0$ 与 $y$ 轴相交于 $A,B$ 两点。求：
  \begin{enumerate}[label=(\alph*),leftmargin=2em]
    \item 圆心及半径；\hfill (2\%)
    \item $AB$ 的长。\hfill (2\%)
  \end{enumerate}

\item 七个小矮人的身高（cm）为：121，116，112，124，127，115，118。
  \begin{enumerate}[label=(\alph*),leftmargin=2em]
    \item 小红帽加入这群小矮人后，平均身高没有改变。求小红帽的身高。\hfill (2\%)
    \item 小木偶加入这群小矮人后，身高的中位数增加了 $1\text{ cm}$。求小木偶的身高。\hfill (2\%)
  \end{enumerate}

\item 已知 $f(x)=x^{-\frac13}$，求 $f'(64)$，并用近似计算求 $\sqrt[3]{\dfrac1{66}}$ 的近似值，答案以分数表示。\hfill (4\%)
\end{enumerate}

\vfill
\begin{flushright}图1：原卷附图\end{flushright}

\newpage

% Page 13
{\bfseries 乙组\quad 选答题\,(40\%)}\par
（本组7题选答4题，但不能超过4题。）

\begin{enumerate}[leftmargin=2.2em,label=\arabic*.,start=6]
\item \begin{enumerate}[label=(\alph*),leftmargin=2em]
  \item 解不等式 $3^x<x+4$。\hfill (4\%)
  \item 解方程式 $4^x-3^x=x$。\hfill (3\%)
  \item 一银行的年利率为 $5\%$。若以复利计算，每半年计息一次，至少要经过多少年，存款才是本金的两倍？\hfill (3\%)
\end{enumerate}

\item \begin{enumerate}[label=(\alph*),leftmargin=2em]
  \item 已知 $26x-41\equiv A(x-3)(x-2)+B(x-3)+C(x-2)$，求 $A,B,C$ 的值。\hfill (3\%)
  \item 解方程式 $\log_5 4+\log_5 3+\log_5(\log x)=0$。\hfill (3\%)
  \item 已知等差数列 $553,546,539,\ldots$。若此数列的首 $n$ 项之和为负值，求 $n$ 的最小值。\hfill (4\%)
\end{enumerate}

\item \begin{enumerate}[label=(\alph*),leftmargin=2em]
  \item 在 $\triangle ABC$ 中，$c=2$，$a=1$，$\cos B=\dfrac35$，求 $\triangle ABC$ 的面积。\hfill (3\%)
  \item 已知 $\tan x=4$，不许使用计算机，求 $\dfrac{\sin x+\cos x}{\sin x-\cos x}$ 的值。\hfill (2\%)
  \item 将 $3\cos\theta-3\sin\theta$ 化为 $R\cos(\theta+\alpha)$ 的形式，其中 $R>0$，$0^\circ<\alpha<90^\circ$。据此解方程式 $3\cos\theta-3\sin\theta=2$，$0^\circ\le\theta\le360^\circ$（答案准确至两位小数）。\hfill (5\%)
\end{enumerate}

\item \begin{enumerate}[label=(\alph*),leftmargin=2em]
  \item 已知平行四边形 $ABCD$ 的三个顶点 $A(-1,2)$，$B(3,1)$ 及 $C(2,4)$。不许用图解法，求：
    \begin{enumerate}[label=(\roman*),leftmargin=2em]
      \item 点 $D$ 的坐标；\hfill (2\%)
      \item 平行四边形 $ABCD$ 的面积；\hfill (2\%)
      \item 点 $B$ 到直线 $AC$ 的距离。\hfill (2\%)
    \end{enumerate}
  \item 若直线 $y=kx$ 是圆 $(x-1)^2+(y-2)^2=4$ 的切线，求 $k$ 的值。\hfill (4\%)
\end{enumerate}
\end{enumerate}

\newpage

% Page 14
\begin{enumerate}[leftmargin=2.2em,label=\arabic*.,start=10]
\item \begin{enumerate}[label=(\alph*),leftmargin=2em]
  \item 有几种方法将十二名学生分成三个小组且每组有四名学生？\hfill (3\%)
  \item 一项投资计划的盈利胥视市场的情况而定，其关系以及对市场的预测如表1所示。求这项投资计划盈利的期望值。\hfill (3\%)
\end{enumerate}

\begin{center}
\begin{tabular}{cccc}
\toprule
市场情况 & 市场良好 & 市场稳定 & 市场低迷 \\
\midrule
盈利 & RM 200,000 & RM 50,000 & $-$RM 350,000 \\
概率 & 0.35 & 0.45 & 0.20 \\
\bottomrule
\end{tabular}
\par 表1
\end{center}

\begin{enumerate}[label=(\alph*),leftmargin=2em,start=3]
  \item 表2显示五名学生每天上网的平均时间与他们的学期成绩总平均：求这些学生每天上网的平均时间（小时）与学期成绩总平均的相关系数。\hfill (4\%)
\end{enumerate}

\begin{center}
\begin{tabular}{cccccc}
\toprule
学生 & 甲 & 乙 & 丙 & 丁 & 戊 \\
\midrule
平均上网时间（小时） & 1.5 & 2 & 0.5 & 1 & 4 \\
学期成绩总平均 & 62.3 & 60.7 & 72.5 & 73.4 & 55.5 \\
\bottomrule
\end{tabular}
\par 表2
\end{center}

\item \begin{enumerate}[label=(\alph*),leftmargin=2em]
  \item 已知函数 $f(x)=x^2-3xe^x$，求 $f$ 为增函数的区间。\hfill (4\%)
  \item 已知一圆柱体的底圆半径为 $r\text{ cm}$，高为 $h\text{ cm}$ 且其体积为 $250\pi\text{ cm}^3$。
  \begin{enumerate}[label=(\roman*),leftmargin=2em]
    \item 证明此圆柱体的表面积 $A=2\pi r^2+\dfrac{500\pi}{r}\text{ cm}^2$。\hfill (2\%)
    \item 求 $r$ 使得圆柱体的表面积为最小。\hfill (4\%)
  \end{enumerate}
\end{enumerate}

\item \begin{enumerate}[label=(\alph*),leftmargin=2em]
  \item 计算 $\displaystyle \lim_{x\to\infty}\frac{4x^3-12x^2}{2x^3+3x-2}$。\hfill (2\%)
  \item 求 $\displaystyle \int \frac{dx}{x^2+4x+4}$。\hfill (2\%)
  \item 图2中的阴影区域是由曲线 $y=4-x^2$ 与直线 $x+2y=4$ 所围成。求此区域绕 $x$ 轴旋转 $360^\circ$ 所形成的旋转体体积。\hfill (6\%)
\end{enumerate}

\vfill
\begin{flushright}图2：原卷附图\end{flushright}

\end{enumerate}

\end{document}